\chapter{Squared-TT Recursion and Fixed-Branch Likelihood Construction}
\label{ch:bf-highdim-squared-tt-fixed-branch-likelihoods}

\section{Purpose And Scope}
\label{sec:bf-hd-squared-tt-purpose}

The previous chapter fixed the retained-object front half of the non-Gaussian
lane: the exact adjacent-state target, the TT approximation toolkit, the
coordinate systems, the KR/preconditioning map logic, and the meaning of the
retained marginal \(a_t\) and normalized carried filter \(\widehat p_t\).
Those objects answer what the method is trying to represent.  The present chapter
answers the operational question that follows: once the branch has committed to a
finite-dimensional TT/KR construction, what exact approximate likelihood object
is produced, what must be saved for the next step, and which numerical fields are
part of the scalar identity rather than merely transient implementation detail?

This chapter is therefore the operational middle of the p50 lane.  It is not yet
the full derivative chapter.  It builds the squared-TT recursion, the retained
marginal and normalizer, the fixed least-squares branch specialization, and the
stored-field contract that later allows the shared same-scalar gradient chapter
to differentiate one declared scalar rather than an adaptively changing routine.

\section{Squared-TT Filtering Object And Recursive Closure}
\label{sec:bf-hd-squared-tt-closure}

The retained-object lane must answer two questions at the same time:
\begin{equation}
  \text{How do we represent } q_t?
  \qquad
  \text{How do we integrate the represented object?}
  \label{eq:bf-hd-squared-tt-two-questions}
\end{equation}
The exact adjacent-state target is a many-variable function rather than a
finite-dimensional Gaussian summary.  In the TT/KR lane, the target is first
moved into reference coordinates and then approximated by a nonnegative
low-rank object, typically through a square or squared-density representation.

Write the represented adjacent-state object in retained coordinates as
\begin{equation}
  \widehat q_t(z_t,z_{t-1}),
  \label{eq:bf-hd-squared-tt-adjacent-object}
\end{equation}
and define the retained numerator and normalizer by
\begin{equation}
  a_t(z_t)=\int \widehat q_t(z_t,z_{t-1})\dd z_{t-1},
  \qquad
  \widehat Z_t=\int a_t(z_t)\dd z_t.
  \label{eq:bf-hd-squared-tt-retained-numerator}
\end{equation}
The carried filter is then
\begin{equation}
  \widehat p_t(z_t)=\frac{a_t(z_t)}{\widehat Z_t}.
  \label{eq:bf-hd-squared-tt-carried-filter}
\end{equation}
This is the recursive closure the lane must preserve: a represented adjacent
object, a retained numerator, a normalizer, and a normalized carried filter that
can be queried at the next time step.  Saving only samples, only a low-rank core
without normalization metadata, or only \(\widehat Z_t\) is not enough to close
the same branch recursively.

\section{What Must Be Saved For The Next Time Step}
\label{sec:bf-hd-squared-tt-stored-fields}

The next subtlety is what must survive from one time step to the next.  If too
little is stored, the next target cannot be reconstructed on the same branch.  If
the wrong object is stored, the implementation may look internally coherent while
no longer evaluating the same scalar as the fixed branch.

For a fixed-branch likelihood implementation, the forward object that must be
available at time \(t+1\) is not merely a cloud or a normalizer.  It is the
retained filter evaluator and the ingredients required to differentiate that
retained filter on the same branch.  A useful stored-field contract is
\begin{equation}
  \mathcal F_t=
  \{C_{t,k},\ M_{t,>k},\ M_{t,<k},\ \widehat Z_t,\ a_t,\ \widehat p_t,
    \Psi_t,\ \lambda_t,\ \tau_t\}_{k=1}^{D},
  \label{eq:bf-hd-squared-tt-stored-forward}
\end{equation}
with the derivative chapter later augmenting it by the corresponding dotted
objects.

For the derivative-aware fixed branch, the stronger stored-field contract is
\begin{equation}
  \dot{\mathcal F}_t=
  \{\dot C_{t,k}^{(i)},\ \dot M_{t,>k}^{(i)},\ \dot M_{t,<k}^{(i)},\
    \dot{\widehat Z}_t^{(i)},\ \dot a_t^{(i)},\ \dot{\widehat p}_t^{(i)}\}_{k=1}^{D},
  \label{eq:bf-hd-squared-tt-stored-derivative}
\end{equation}
so that the next transformed target may be evaluated together with its product-rule
sensitivity on the same branch.  The carried-filter derivative is the quotient
object
\begin{equation}
  \partial_{\beta_i}\widehat p_t(x_t;\beta)
  =
  \frac{\partial_{\beta_i}a_t(x_t;\beta)}{\widehat Z_t(\beta)}
  -
  \frac{a_t(x_t;\beta)\partial_{\beta_i}\widehat Z_t(\beta)}{\widehat Z_t(\beta)^2},
  \label{eq:bf-hd-squared-tt-carried-derivative}
\end{equation}
which is precisely the object needed by the next-step fixed-branch derivative
pass.

The point is conceptual before it is algorithmic: the next-step target depends
on evaluating the previous carried filter at new fitting points.  Therefore the
retained-object lane must store the represented numerator/filter object in a form
that the next-step evaluator can query without rebuilding an adaptive branch.
This is what makes the fixed-branch recursion meaningfully recursive rather than a
sequence of disconnected local fits.

\section{Fixed-Branch Likelihood Construction}
\label{sec:bf-hd-squared-tt-fixed-branch-likelihood}

The adaptive Zhao--Cui algorithm contains discrete choices---ranks, fitting
points, sweep decisions, domain maps, and stabilization choices---that should not
be treated as globally differentiable with respect to the scientific parameter.
The fixed-branch specialization freezes those discrete choices first, then asks
what scalar is defined when the corresponding numerical equations are recomputed
on that same frozen structure.

The fixed branch is therefore more than a density formula.  At time \(t\), the
finite-dimensional branch data may be summarized as
\begin{equation}
  \mathcal B_t=
  \{\Omega_{t,k},\Psi_{t,k},\psi_{t,k,\ell},C_{t,k},R_{t,k},B_{t,k},
    \lambda_{t,k},\tau_t,c_t,\text{mass contractions},\text{map data}\}.
  \label{eq:bf-hd-squared-tt-branch-data}
\end{equation}
Changing domains, bases, rank tuples, TT cores, reference densities, defensive
mass conventions, or log-scale shifts changes the approximate likelihood.  Those
objects are part of the scalar identity, not optional implementation metadata.

\subsection{What A Core Object Looks Like}
\label{sec:bf-hd-squared-tt-core-object}

For coordinate \(k\), store
\begin{equation}
  C_k\in\R^{R_{k-1}\times p_k\times R_k}.
  \label{eq:bf-hd-squared-tt-core-shape}
\end{equation}
Given a scalar coordinate value \(r_k\), first compute the basis vector
\begin{equation}
  b_k(r_k)=
  \bigl(\psi_{k,0}(r_k),\ldots,\psi_{k,p_k-1}(r_k)\bigr)^\top.
  \label{eq:bf-hd-squared-tt-basis-vector}
\end{equation}
Then form the matrix slice
\begin{equation}
  H_k(r_k)[a,b]
  =
  \sum_{\ell=0}^{p_k-1}C_k[a,\ell,b]b_k(r_k)[\ell].
  \label{eq:bf-hd-squared-tt-core-eval}
\end{equation}
Evaluation of the TT is the loop
\begin{equation}
  v_0=1,
  \qquad
  v_k=v_{k-1}H_k(r_k),\quad k=1,\ldots,D,
  \qquad
  \widehat h(r)=v_D.
  \label{eq:bf-hd-squared-tt-tt-eval}
\end{equation}
The square-density object additionally stores mass contractions
\begin{equation}
  M_{>k}
  =
  \int H_{k+1}(r_{k+1})\cdots H_D(r_D)
  H_D(r_D)^\top\cdots H_{k+1}(r_{k+1})^\top
  \dd r_{k+1:D},
  \label{eq:bf-hd-squared-tt-mass-contraction}
\end{equation}
with analogous left-side contractions for reverse-order marginals.  These are
not decorative arrays; they are the finite-dimensional objects through which
marginals, normalizers, and later derivatives are actually computed.

The resulting fixed-branch scalar has the form
\begin{equation}
  \widehat\ell_T^{\rm TT/KR}(\beta;B)
  =\sum_{t=1}^T \log \widehat Z_t(\beta;B_t),
  \label{eq:bf-hd-squared-tt-scalar}
\end{equation}
where the adaptive selection procedure is no longer being differentiated.  What
is being differentiated later is the scalar induced by recomputing the fitted
numerical objects on this saved branch.

\section{Consolidated Fixed Least-Squares Formulation}
\label{sec:bf-hd-squared-tt-fixed-ls}

The implementable fixed-branch route in BayesFilter should be stated narrowly.
It is not adaptive TT-cross.  It is a fixed least-squares specialization with
fixed points, fixed ranks, fixed sweep order, and deterministic stabilization.
In symbols,
\begin{equation}
  \text{fixed-branch TT likelihood route}
  =
  \text{fixed points}
  +
  \text{fixed ranks}
  +
  \text{fixed sweep order}
  +
  \text{deterministic least squares / stabilization}.
  \label{eq:bf-hd-squared-tt-fixed-ls-summary}
\end{equation}
This narrower statement matters because it isolates the branch whose value and
later derivative can be audited.  If the implementation silently switches pivot
sets, changes rank caps, or changes stabilization rules under perturbation, then
it has left the fixed-branch scalar described here.

A representative one-step input/output contract is:
\begin{align}
  \mathcal I_t &=
  \{y_t,\beta,\widehat p_{t-1},f_\beta,g_\beta,\Psi_t,Z_{\rm fit}^{(t)},W^{(t)},
    \text{basis},\text{ranks},\text{sweeps},\rho_t,c_t,\tau_t\},
  \label{eq:bf-hd-squared-tt-inputs}\\
  \mathcal O_t &=
  \{C_{t,k},M_{t,>k},M_{t,<k},\widehat Z_t,a_t,\widehat p_t\}_{k=1}^D.
  \label{eq:bf-hd-squared-tt-outputs}
\end{align}
The important point is not the exact symbol list but the discipline: the branch
must state what is input, what is output, and which outputs are required to
close the next-step recursion on the same branch.

A concrete shape contract helps keep this branch auditable.  One useful table is
\begin{equation}
\begin{array}{c|c}
\text{object} & \text{shape}\\
\hline
C_k & R_{k-1}\times p_k\times R_k\\
L_{j,k} & 1\times R_{k-1}\\
R_{j,k} & R_k\times 1\\
A_k & n_{\rm fit}\times (R_{k-1}p_kR_k)\\
g_k=\operatorname{vec}(C_k),\ d_k & R_{k-1}p_kR_k\\
N_k & (R_{k-1}p_kR_k)\times (R_{k-1}p_kR_k)
\end{array}
  \label{eq:bf-hd-squared-tt-shape-table}
\end{equation}
with vectorization convention
\begin{equation}
  \iota_k(a,\ell,b)=
  (a-1)p_kR_k+\ell R_k+b,
  \qquad
  g_k[\iota_k(a,\ell,b)]=C_k[a,\ell,b].
  \label{eq:bf-hd-squared-tt-flattening}
\end{equation}
This is the level at which the later derivative chapter can say, concretely, what
it means to recompute the same least-squares solve rather than copying a stale
core from the base parameter.

\subsection{One Concrete Branch, Fully Instantiated}
\label{sec:bf-hd-squared-tt-concrete-branch}

A branch record is easier to audit when one concrete instantiation is stated.
For a declared run, that record should pin down:
\begin{enumerate}[leftmargin=0.28in]
  \item coordinate order and domain policy,
  \item basis family and basis degree in each coordinate,
  \item deterministic fitting points and weights,
  \item rank vector and deterministic rank ladder before branch freezing,
  \item sweep order and number of sweeps,
  \item defensive reference density, defensive mass, and any frozen shift
  convention,
  \item the boxed convention that the fit is performed in source scale while the
  reported scalar returns to the unshifted evidence scale.
\end{enumerate}
This list is not busywork.  If one of these items changes under perturbation,
then the implementation has left the fixed-branch scalar even if all later
notation still looks similar.

A useful first mathematical instantiation is intentionally narrow.  For a scalar
state, use the two-block coordinate order
\begin{equation}
  r_t=(x_t,x_{t-1})\in\R^2.
  \label{eq:bf-hd-squared-tt-concrete-order}
\end{equation}
For vector states, concatenate coordinates as
\begin{equation}
  r_t=(x_{t,1},\ldots,x_{t,m_x},x_{t-1,1},\ldots,x_{t-1,m_x}).
  \label{eq:bf-hd-squared-tt-concrete-vector-order}
\end{equation}
Choose a fixed interval for each coordinate,
\begin{equation}
  r_{t,k}\in[\ell_{t,k},u_{t,k}],
  \qquad
  r_{t,k}=a_{t,k}+h_{t,k}z_k,
  \qquad
  a_{t,k}=\frac{u_{t,k}+\ell_{t,k}}2,
  \qquad
  h_{t,k}=\frac{u_{t,k}-\ell_{t,k}}2,
  \label{eq:bf-hd-squared-tt-concrete-domain}
\end{equation}
so the Jacobian is constant,
\begin{equation}
  J_t=\prod_{k=1}^Dh_{t,k}.
  \label{eq:bf-hd-squared-tt-concrete-jacobian}
\end{equation}
The branch stores \(\ell_{t,k},u_{t,k},a_{t,k},h_{t,k}\), and these remain fixed
through same-branch finite differences.

A deterministic domain-selection policy should be declared before branch
freezing.  Let a pilot location and scale be
\begin{equation}
  \mu_{t,k}^{\rm pilot},
  \qquad
  s_{t,k}^{\rm pilot}>0,
  \label{eq:bf-hd-squared-tt-pilot-domain}
\end{equation}
and set
\begin{equation}
  [\ell_{t,k},u_{t,k}]
  =
  [\mu_{t,k}^{\rm pilot}-\gamma s_{t,k}^{\rm pilot},
   \mu_{t,k}^{\rm pilot}+\gamma s_{t,k}^{\rm pilot}],
  \label{eq:bf-hd-squared-tt-domain-policy}
\end{equation}
with fixed expansion factor \(\gamma\).  In the current fixed-variant fitting
story, these pilot moments may come from a Gaussian scout such as a UKF warm
start.  The scout provides a local frame and an initialization region; it is not
a certificate that the represented density is already correct.  The
out-of-domain diagnostic at fitting or proposal points is
\begin{equation}
  \delta_{\rm out}=
  \frac1N\sum_{j=1}^N \mathbf{1}\{r_{t,k}^{(j)}\notin[\ell_{t,k},u_{t,k}]\ \text{for some }k\},
  \label{eq:bf-hd-squared-tt-domain-diagnostic}
\end{equation}
and the branch is accepted only if
\begin{equation}
  \delta_{\rm out}\le\delta_{\max}.
  \label{eq:bf-hd-squared-tt-domain-accept}
\end{equation}
If this fails, the branch must be rebuilt before differentiation rather than
expanded during a same-branch derivative check.

Use one declared basis family and fixed point design.  A useful first choice is
normalized Legendre basis of common degree \(p\) in every coordinate together
with deterministic Halton-type fitting points and equal weights.  The rank vector
is then selected by a small deterministic ladder before branch freezing, for
example
\begin{equation}
  \mathcal R_{\rm ladder}=
  \{(1,\ldots,1),(2,\ldots,2),(4,\ldots,4),(6,\ldots,6),(8,\ldots,8)\},
  \label{eq:bf-hd-squared-tt-rank-ladder}
\end{equation}
with the first candidate passing the declared residual, conditioning, and
defensive-mass tests becoming the frozen branch rank.  If no candidate passes,
the correct outcome is branch failure, not an unrecorded adaptive rank change.

The defensive reference and shift convention should also be fixed explicitly.  A
minimal deterministic choice is uniform reference density on \([-1,1]^D\), a
source-scale floor \(\epsilon_{\rm floor}>0\), defensive mass
\(\tau_t=2^D\epsilon_{\rm floor}\), and shift \(c_t\) computed once at the base
parameter and then reused for same-branch finite differences.  This is the
meaning of the chapter's boxed convention that the fit is carried out in source
scale while the reported scalar returns to the unshifted evidence scale.

\subsection{Square-Root Fitting Object Versus Density Audit Object}
\label{sec:bf-hd-squared-tt-density-audit}

The fixed-variant fitting lane needs one further distinction.  The TT sweep may
be written in square-root form, but the scientific object being evaluated is a
density.  Write the represented nonnegative object as
\begin{equation}
  q_\theta(z)=h_\theta(z)^2+\tau q_0(z),
  \qquad
  Z_\theta=\int q_\theta(z)\dd z,
  \qquad
  p_\theta(z)=\frac{q_\theta(z)}{Z_\theta},
  \label{eq:bf-hd-squared-tt-density-object}
\end{equation}
where \(h_\theta\) is the trainable square-root TT field, \(q_0\) is the fixed
defensive reference density, and \(\tau>0\) is the defensive mixture weight.
If the target square-root field is \(s(z)\), then the target density is
\begin{equation}
  p_\star(z)\propto s(z)^2.
  \label{eq:bf-hd-squared-tt-target-density}
\end{equation}
Therefore a raw square-root residual,
\begin{equation}
  h_\theta(z_i)\approx s(z_i),
  \label{eq:bf-hd-squared-tt-square-root-residual}
\end{equation}
is not yet the primary heldout criterion.  The reported density depends on the
square, the defensive component, and the normalizer \(Z_\theta\).  Sign, scale,
and local amplitude errors in \(h_\theta\) need not map one-for-one to density
quality after normalization.  Square-root residuals are therefore fitting
diagnostics, not promotion criteria by themselves.

For a weighted batch
\begin{equation}
  B=\{(z_i,w_i,s_i)\}_{i=1}^n,
  \label{eq:bf-hd-squared-tt-heldout-batch}
\end{equation}
the target-only empirical measure is
\begin{equation}
  \alpha_i^{\rm hold}
  =
  \frac{w_i s_i^2}{\sum_j w_j s_j^2},
  \label{eq:bf-hd-squared-tt-heldout-alpha}
\end{equation}
and the corresponding density cross-entropy is
\begin{equation}
  \widehat{\mathcal L}_{\rm hold}(\theta)
  =
  -\sum_i \alpha_i^{\rm hold}\log q_\theta(z_i)
  +\log Z_\theta.
  \label{eq:bf-hd-squared-tt-heldout-cross-entropy}
\end{equation}
This target-only measure is the corrected heldout audit object for the
fixed-variant lane.  The historical helper
\begin{equation}
  \alpha_i^{\rm old}
  \propto
  w_i\bigl(s_i^2+\tau q_0(z_i)\bigr)
  \label{eq:bf-hd-squared-tt-old-helper-alpha}
\end{equation}
mixes the defensive reference density into the empirical target.  That helper may
still be useful as a boundary comparator or internal training-support convention,
but it is not the heldout target measure and should not be reported as if it
were the scientific density criterion.

This correction also changes how fitting evidence should be narrated.  A UKF
warm start can supply local geometry and a mini-batch training run can improve
the represented density over that geometry, but success is judged by the
target-only heldout cross-entropy in
\eqref{eq:bf-hd-squared-tt-heldout-cross-entropy}, not by the old
\(\tau q_0\)-smoothed helper and not by square-root residuals alone.

\section{Initialization And Sweep Order}
\label{sec:bf-hd-squared-tt-init-sweeps}

The branch must also save the deterministic initialization and sweep rule.  A
constant-channel initialization is a useful mechanical baseline, but it should
not be confused with the preferred fixed-variant fitting narrative once a
deterministic geometric warm start is available.  When a Gaussian scout such as
a UKF provides the pilot frame, the current mathematical story is that the
branch record should freeze a warm start derived from that scout rather than
revert to random or purely constant initialization.  Let one bidirectional sweep
mean updates in the order
\begin{equation}
  1,2,\ldots,D,D-1,\ldots,1.
  \label{eq:bf-hd-squared-tt-sweep-order}
\end{equation}
Fix the number of bidirectional sweeps \(S\).  The exact deterministic
initialization rule matters less than saving it in the branch record, because
changing initialization changes the fixed fitting map.

\subsection{Forward Core Update}
\label{sec:bf-hd-squared-tt-forward-core-update}

At fitting point \(z^{(j)}\), define environments for core \(k\) by
\begin{equation}
  L_{j,k}=H_1(z_1^{(j)})\cdots H_{k-1}(z_{k-1}^{(j)}),
  \qquad
  R_{j,k}=H_{k+1}(z_{k+1}^{(j)})\cdots H_D(z_D^{(j)}).
  \label{eq:bf-hd-squared-tt-forward-environments}
\end{equation}
For endpoint ranks this means \(L_{j,k}\in\R^{1\times R_{k-1}}\) and
\(R_{j,k}\in\R^{R_k\times1}\).  The design row for coefficient
\((a,\ell,b)\) is
\begin{equation}
  A_k[j,(a,\ell,b)]
  =
  L_{j,k}[1,a]\,b_k(z_k^{(j)})[\ell]R_{j,k}[b,1].
  \label{eq:bf-hd-squared-tt-forward-design-row}
\end{equation}
The local ridge system is
\begin{equation}
  N_k g_k=d_k,
  \qquad
  N_k=A_k^\top W A_k+\rho I,
  \qquad
  d_k=A_k^\top Wy.
  \label{eq:bf-hd-squared-tt-forward-ridge-system}
\end{equation}
After solving \eqref{eq:bf-hd-squared-tt-forward-ridge-system}, reshape \(g_k\)
back into the core by
\begin{equation}
  C_k[a,\ell,b]=g_k[(a,\ell,b)].
  \label{eq:bf-hd-squared-tt-forward-core-reshape}
\end{equation}
This is the forward fitting rule whose identity later same-branch finite
differences must preserve.

\subsection{Derivative Core Update}
\label{sec:bf-hd-squared-tt-derivative-core-update}

The derivative environments are obtained by differentiating the matrix products:
\begin{equation}
  \dot L_{j,k}
  =
  \sum_{r=1}^{k-1}H_1\cdots H_{r-1}\dot H_r H_{r+1}\cdots H_{k-1},
  \label{eq:bf-hd-squared-tt-dot-L}
\end{equation}
\begin{equation}
  \dot R_{j,k}
  =
  \sum_{r=k+1}^{D}H_{k+1}\cdots H_{r-1}\dot H_r H_{r+1}\cdots H_D.
  \label{eq:bf-hd-squared-tt-dot-R}
\end{equation}
Because basis functions and fitting points are frozen, the derivative design row
contains only environment terms:
\begin{equation}
\begin{aligned}
  \dot A_k[j,(a,\ell,b)]
  &=
  \dot L_{j,k}[1,a]b_k(z_k^{(j)})[\ell]R_{j,k}[b,1]\\
  &\quad+
  L_{j,k}[1,a]b_k(z_k^{(j)})[\ell]\dot R_{j,k}[b,1].
\end{aligned}
  \label{eq:bf-hd-squared-tt-dot-design-row}
\end{equation}
The derivative normal equations are
\begin{equation}
  \dot N_k=\dot A_k^\top W A_k+A_k^\top W\dot A_k,
  \qquad
  \dot d_k=\dot A_k^\top Wy+A_k^\top W\dot y,
  \label{eq:bf-hd-squared-tt-dot-normal-equations}
\end{equation}
and differentiating \(N_kg_k=d_k\) gives
\begin{equation}
  N_k\dot g_k=\dot d_k-\dot N_k g_k.
  \label{eq:bf-hd-squared-tt-dot-ridge-system}
\end{equation}
After solving \eqref{eq:bf-hd-squared-tt-dot-ridge-system}, reshape
\begin{equation}
  \dot C_k[a,\ell,b]=\dot g_k[(a,\ell,b)].
  \label{eq:bf-hd-squared-tt-dot-core-reshape}
\end{equation}
This is the derivative fitting rule.  A later implementation must use this same
stored computational path rather than inventing a separate approximate derivative
solve.

\subsection{Alternating-Sweep Recomputations}
\label{sec:bf-hd-squared-tt-sweep-recomputations}

The symbols \(L_{j,k}\), \(R_{j,k}\), \(A_k\), \(N_k\), and \(d_k\) are not
static arrays during an alternating sweep.  They are recomputed from the current
cores at the moment a core is updated.  Let
\begin{equation}
  H_{j,k}=H_k(z_k^{(j)};C_k),
  \qquad
  \dot H_{j,k}=H_k(z_k^{(j)};\dot C_k),
  \label{eq:bf-hd-squared-tt-sweep-H}
\end{equation}
where basis values and fitting points are fixed and only the coefficients change.
For a current core collection \(C^{\rm cur}\), define prefix and suffix products
by
\begin{equation}
  L_{j,1}^{\rm cur}=1,
  \qquad
  L_{j,k}^{\rm cur}=L_{j,k-1}^{\rm cur}H_{j,k-1}^{\rm cur},
  \qquad k=2,\ldots,D,
  \label{eq:bf-hd-squared-tt-sweep-prefix}
\end{equation}
and
\begin{equation}
  R_{j,D}^{\rm cur}=1,
  \qquad
  R_{j,k}^{\rm cur}=H_{j,k+1}^{\rm cur}R_{j,k+1}^{\rm cur},
  \qquad k=D-1,\ldots,1.
  \label{eq:bf-hd-squared-tt-sweep-suffix}
\end{equation}
The dotted prefix and suffix products follow the same recursions with one
product-rule term added.  At a core update \(k\), build the current design
matrix and normal equations from the current environments, solve for the new
core, then build the dotted objects and solve the derivative system.  After
updating one core, every cached prefix with index \(\ell>k\) and every cached
suffix with index \(\ell<k\) is stale and must be recomputed before use.

This stale-cache rule is part of the branch definition.  It is what makes the
phrase “the same fitting rule” operational rather than rhetorical.

\subsection{Stopping Rule And Failure Exits}
\label{sec:bf-hd-squared-tt-stopping-failures}

After each full sweep compute
\begin{equation}
  \varepsilon_s=
  \frac{\|A(C^{(s)})-y\|_W}{\max\{\|y\|_W,\epsilon_{\rm abs}\}},
  \qquad
  \Delta_s=
  \frac{\|C^{(s)}-C^{(s-1)}\|}{\max\{\|C^{(s-1)}\|,\epsilon_{\rm abs}\}},
  \label{eq:bf-hd-squared-tt-stopping-metrics}
\end{equation}
The sweep stops successfully only if
\begin{equation}
  \varepsilon_s\le\epsilon_{\rm fit}
  \quad\text{and}\quad
  \Delta_s\le\epsilon_{\rm sweep}.
  \label{eq:bf-hd-squared-tt-stopping-success}
\end{equation}
If this does not happen by \(S\) sweeps, the branch reports sweep failure rather
than silently accepting a different scalar.  The full accepted/failure path must
therefore distinguish at least the following exits:
\begin{center}
\begin{tabular}{p{0.28\linewidth}p{0.34\linewidth}p{0.28\linewidth}}
\toprule
Failure condition & Mathematical test & Report field\\
\midrule
Domain miss & \(\delta_{\rm out}>\delta_{\max}\) & domain failure\\
Rank ladder fail & no \(R\) satisfies residual/conditioning/floor tests & rank failure\\
Ill-conditioned solve & \(\operatorname{cond}(N_k)>\kappa_{\max}\) & solve failure\\
Sweep stagnation & \(\varepsilon_s>\epsilon_{\rm fit}\) or \(\Delta_s>\epsilon_{\rm sweep}\) after \(S\) sweeps & sweep failure\\
Nonpositive target for smooth derivative & \(\min_j\widetilde q_{t,j}\le0\) without declared floor & target failure\\
Mass failure & \(\widehat Z_t\le0\) or not finite & normalizer failure\\
Finite-difference branch mismatch & frozen fields or accepted/failure path differ across \(\beta-h,\beta,\beta+h\) & branch mismatch\\
\bottomrule
\end{tabular}
\end{center}
A failed branch is not a successful derivative with a warning attached; it is a
declared failure of the fixed scalar.  This accepted/failure vocabulary is what
staged \(ch37\) later consumes when it says a branch-valid FD row requires the
same realized acceptance/failure path on both sides of the perturbation.

\section{One Observation Step Through The Stored Objects}
\label{sec:bf-hd-squared-tt-one-step}

A reader should be able to narrate one time step in the following order:
\begin{enumerate}[leftmargin=0.28in]
  \item evaluate the previous carried filter on the current branch’s fitting
  points,
  \item combine it with the transition and observation densities to form the
  adjacent-state target in the declared coordinates,
  \item fit the squared-TT / low-rank object on the fixed branch,
  \item contract out the previous-state block to form \(a_t\),
  \item compute \(\widehat Z_t\) and normalize to \(\widehat p_t\),
  \item save the branch fields needed for the next-step forward evaluation and
  later same-branch derivative evaluation.
\end{enumerate}

The dependency flow should be explicit enough to implement:
\begin{equation}
\begin{array}{rcl}
\text{reference points} &\longrightarrow& z^{(j)}\ \text{from the fixed point design},\\
\text{physical points} &\longrightarrow& r^{(j)}=\Psi_t(z^{(j)}),\\
\text{target values} &\longrightarrow&
\widetilde q_{t,j}=
\widehat p_{t-1}(x_{t-1}^{(j)};\beta)
 f_\beta(x_t^{(j)}\mid x_{t-1}^{(j)})
 g_\beta(y_t\mid x_t^{(j)})J_t,\\
\text{square-root fit values} &\longrightarrow&
 y_j=e^{c_t/2}\sqrt{\widetilde q_{t,j}},\\
\text{initial cores} &\longrightarrow&
 \text{deterministic constant-channel initialization},\\
\text{fixed sweep objects} &\longrightarrow&
 L_{j,k},R_{j,k},A_k,N_k,d_k,C_k,\\
\text{mass objects} &\longrightarrow&
 M_{>k},M_{<k},R_t,\widehat Z_t,\\
\text{carried objects} &\longrightarrow&
 a_t,\widehat p_t,\\
\text{saved objects for }t+1 &\longrightarrow&
 \mathcal F_t\ \text{and, when needed, }\dot{\mathcal F}_t.
\end{array}
  \label{eq:bf-hd-squared-tt-forward-flow}
\end{equation}
The forward environments are recomputed from the current cores during each sweep;
they are not static arrays.  For fitting point \(z^{(j)}\), the environments are
\begin{equation}
  L_{j,k}=H_1(z_1^{(j)})\cdots H_{k-1}(z_{k-1}^{(j)}),
  \qquad
  R_{j,k}=H_{k+1}(z_{k+1}^{(j)})\cdots H_D(z_D^{(j)}),
  \label{eq:bf-hd-squared-tt-environments}
\end{equation}
and the local ridge system is built from those current environments.  After a
core update, cached prefixes and suffixes on the wrong side of that core are
stale and must be recomputed before use.  That bookkeeping is part of the fixed
fitting map.

The same principle governs next-step closure.  The next transformed target is not
formed from a vague “previous TT object,” but from the saved carried evaluator:
\begin{equation}
  \widetilde q_{t+1}(z;\beta)=
  \widehat p_t(x_t(z);\beta)
  f_\beta(x_{t+1}(z)\mid x_t(z))
  g_\beta(y_{t+1}\mid x_{t+1}(z))J_{t+1}(z).
  \label{eq:bf-hd-squared-tt-next-target}
\end{equation}
If the derivative lane is active, the saved dotted object from
\eqref{eq:bf-hd-squared-tt-stored-derivative} is what makes the next-step
product-rule derivative possible on the same branch.  Without that saved object,
a later derivative pass would be differentiating a different computational route.

This is the operational middle that the monograph must carry so that the later
derivative chapter can start from a declared scalar object rather than a vague
reference to “the TT filter.”

\section{Exports To The Same-Scalar Derivative Chapter}
\label{sec:bf-hd-squared-tt-exports}

This chapter exports the following objects forward:
\begin{enumerate}[leftmargin=0.28in]
  \item the represented adjacent-state object and retained marginal recursion,
  \item the stored-field contract for next-step closure,
  \item the fixed-branch scalar \(\widehat\ell_T^{\rm TT/KR}(\beta;B)\),
  \item the fixed-branch least-squares specialization,
  \item the lane-specific statement that domains, points, ranks, maps,
  stabilization rules, and stored-field conventions belong to the scalar identity.
\end{enumerate}
The next chapter may now focus on derivative theory, branch-valid finite
differences, and HMC-facing same-scalar discipline, because the present chapter
has already fixed what scalar object is being differentiated.
